\documentclass[a4paper,10pt]{article}
\usepackage[left=1in,right=1in,top=0.75in,bottom=0.75in]{geometry}
\usepackage{fontawesome5} % For icons
\newcommand{\icontext}[2]{\raisebox{-0.25ex}{#1}\,#2}
\usepackage{hyperref}
\hypersetup{
    colorlinks=true,
    urlcolor=[rgb]{0,0,0.4}, % dark navy instead of bright blue
    linkcolor=black,
    citecolor=black
}
\urlstyle{same}
\usepackage{enumitem}
\setlist[itemize]{leftmargin=1.4em, itemsep=2pt, topsep=2pt, parsep=0pt}
\usepackage[T1]{fontenc}
\usepackage{textcomp}
\usepackage{parskip}
\setlength{\parskip}{2pt}
\usepackage{titlesec}
\usepackage{tabularx}
\usepackage{ragged2e}
\usepackage{siunitx}
\sisetup{group-separator={,}}

\pagestyle{empty}
\setlength{\parindent}{0pt}
\setlength{\parskip}{3pt}

\titleformat{\section}
  {\Large\bfseries\raggedright}
  {}
  {0pt}
  {}
\titlespacing{\section}{0pt}{12pt}{4pt}

\begin{document}

\noindent{\fontsize{18}{20}\selectfont\bfseries Kai Gu\par}
\vspace{2pt}

\faEnvelope\; \href{mailto:kai.gu@studbocconi.it}{kai.gu@studbocconi.it} \quad
\faGlobe\; \href{https://kaigu.github.io}{kaigu.github.io} \quad
\faClock\; Updated: \today

\section*{Education}
\hrule
\vspace{2pt}
\textbf{Bocconi University} \hfill \textbf{Milan, Italy} \\
PhD in Marketing \hfill 2026 -- Present

\vspace{4pt}
\textbf{Bocconi University} \hfill \textbf{Milan, Italy} \\
MSc in Data Science \hfill 2024 -- 2026

\vspace{4pt}
\textbf{Zhejiang University} \hfill \textbf{Hangzhou, China} \\
BEcon in Economics \hfill 2020 -- 2024

\section*{Research Interests}
\hrule
\vspace{2pt}
\textit{Substantive:} Economics of AI, Economics of Digitization \& Platforms, Media Economics \\
\textit{Methodological:} Causal Inference, Computational Social Sciences, Multimodal Data

\section*{Working Papers}
\hrule
\vspace{2pt}
\textbf{Governing AI Content Tools: The Efficiency--Diversity Trade-off in Platform Knowledge Ecosystems}\\
(with \href{https://kaizhu.me/}{Kai Zhu} and \href{https://lorelupo.github.io/}{Lorenzo Lupo})

\smallskip
\begin{quote}\small
Knowledge platforms increasingly deploy AI tools to accelerate content production, yet the ecosystem-level consequences of these deployment decisions remain poorly understood. We conduct a large-scale audit of AI-assisted content creation on Wikipedia, one of the earliest major platform deployments of AI content tools, examining how machine translation (MT) relates to information diversity in multilingual knowledge ecosystems. Using over half a million matched article pairs across ten language editions, we document an efficiency--diversity trade-off: relative to independently written (human-created) target-language articles, MT-assisted articles provide more complete coverage for target-language readers but contribute significantly less novel information to the cross-language knowledge ecosystem. MT-assisted articles exhibit lower semantic divergence from their English source articles and lower information breadth across language pairs and topic categories; the missing content is often culturally anchored---local context, regional perspectives, and cultural details that human editors naturally add. Our findings suggest the need for explicit governance frameworks that acknowledge trade-offs, monitor ecosystem-level outcomes, and design AI deployment strategies that preserve valued properties while capturing efficiency gains.
\end{quote}
\textit{Presented at:} 13th Wiki Workshop (Online), March 2026

\vspace{1em}

\textbf{AI Search and Web Traffic}\\
(with \href{https://qiaonishi.github.io/}{Qiaoni Shi} and \href{https://kaizhu.me/}{Kai Zhu})

\smallskip
\begin{quote}\small
This paper investigates whether AI search restores the referral ecosystem by driving downstream traffic or exacerbates `zero-click' consumption by synthesizing answers within the interface. Leveraging granular clickstream data and a staggered difference-in-differences design, we estimate the causal impact of AI search adoption on users' broader web consumption and compare the magnitude of direct referral traffic against traditional search engines. We further test for heterogeneity to determine if AI intermediaries disproportionately concentrate traffic toward specific domain types and to analyze the `AI divide' across user demographics. Our findings illuminate the long-term viability of the ad-supported internet, providing insights for policymakers and content publishers navigating the evolving competitive dynamics between AI platforms and the open web ecosystem.
\end{quote}
\textit{Presented at:} ISMS Marketing Science Conference, Lisbon, Portugal, June 2026; Bocconi Marketing Research Camp, May 2026

\section*{Conference Presentations}
\hrule
\vspace{2pt}
\textit{ISMS Marketing Science Conference}, Lisbon, Portugal \hfill June 2026 \\
\textit{Bocconi Marketing Research Camp} \hfill May 2026 \\
\textit{13th Wiki Workshop (Online)} \hfill March 2026

\section*{Skills}
\hrule
\vspace{2pt}
\textit{Languages:} English, Mandarin \\
\textit{Programming:} R, Python, \LaTeX, Markdown

\end{document}
